\section{Introduction}
\label{sec:intro}

High-performance numerical software relies on kernels that exploit target-specific instructions and memory behavior. Developing these kernels requires substantial hardware expertise, and moving to a different architecture often entails another round of implementation and tuning. Recent advances in language models have made it possible to automate parts of this work, with agents using hardware descriptions and execution feedback to generate and optimize low-level code~\cite{kernelbench,autocomp}.

We focus on SIMD kernels for CPU vector extensions. Their C intrinsic interfaces expose vector operations within ordinary C programs, allowing agents to modify instruction selection and loop organization directly. These rewrites must account for differences in the underlying execution models. Arm NEON uses fixed-width vectors, whereas RISC-V Vector (RVV) selects the number of active elements at runtime; a port must therefore account for different iteration sizes and tail accesses~\cite{acle,rvv}. Recent work has applied language models to this migration, including translation of NEON kernels to RVV followed by target-specific optimization~\cite{saltyrn,intrintrans}.

Correctness errors can arise in the transformations that improve performance: changing a loop's iteration size affects which elements are loaded, while replacing an intrinsic may change rounding or exceptional-value behavior. These errors can survive tests on ordinary inputs. Existing formal approaches already validate generated code through bounded symbolic execution or compiler-IR reasoning~\cite{saltyrn,llm-vectorizer,alive2}. Establishing equivalence over the full calling domain requires reasoning about every admitted loop iteration and memory access, under a model that retains the relevant ISA differences. Even a one-element maximum can disagree across ISAs if one instruction propagates a NaN and the other selects the numerical operand (Section~\ref{sec:example}). Our objective is to verify the optimized implementation under an explicit calling contract, while preserving the code whose performance motivated the optimization.

This objective requires a semantic model for both the vector operations and the surrounding program. Machine-code and compiler-IR verification provide established ways to express such models~\cite{sail,islaris,alive2,rake}. For intrinsic kernels, the C source already exposes the loops, pointer operations, and calls that must be related. We therefore give the supported intrinsics executable C interpretations and verify the resulting C programs. This keeps the proof obligations close to the source and allows intrinsic definitions and supporting lemmas to be reused across kernels. The interpretations retain architecture-specific behavior; equivalence is established between the complete kernels under their calling contracts.

We present \tool, a framework for semantic equivalence checking of SIMD kernels using reusable C semantic shims. Its design targets NEON, RVV, AVX, and Hexagon HVX. \todo{Report the supported intrinsic subsets.} A common specification determines the required outputs for the logical inputs, while each implementation has its own memory representation and loop invariants. Proving both implementations against this specification relates their results without requiring their loops to advance in lockstep. QCP generates verification conditions for the annotated C programs, which are discharged through automation and Rocq/Coq proofs~\cite{qcp}. These obligations cover the output relation and modeled memory accesses. The guarantee is partial correctness relative to the intrinsic models and calling contracts; Section~\ref{sec:guarantee} states its trust assumptions and scope.

\tool separates the implementation, specification, and proof so that agents can develop proof annotations and auxiliary lemmas without altering the executable algorithm. Experts review the calling contract and any changes suggested during verification. A feasible counterexample calls for repairing the target and checking the revised implementation. A counterexample excluded by the original caller's documented assumptions instead calls for making those assumptions explicit in the specification, with evidence that all legitimate calls remain covered. This separation makes the reason for each change reviewable and allows its performance cost to be measured independently of proof development.

The same verification method supports optimized-kernel validation, translation to another ISA, and cross-checking of independently developed implementations. For translation, an existing NEON or AVX kernel supplies a reference when a target library lacks the operation. For cross-checking, the question is whether implementations of the same operation agree under a shared contract, extending the comparison underlying differential testing to a program equivalence obligation~\cite{mckeeman,symdiff}. A concrete disagreement identifies behavior to examine against the library contract, helping distinguish a defect requiring repair from an intentional numerical difference.

Our evaluation will examine these uses through the 35 final optimized RVV kernels released by SALTyRN, ports of five AVX operations to RVV and HVX, and corresponding implementations in XNNPACK~\cite{final-kernels}. A composed SIMD program, with softmax as a candidate, will test whether the kernel contracts support reasoning across operation boundaries. \todo{Report verification coverage and cost, justified repairs and their performance effects, cross-checking findings, and the composed-program result.} The main measurements are the optimized implementations that can be verified without executable changes, the effort required to extend the semantic models, and the performance retained after necessary repairs.

% Alternative short applications paragraph, retained for later author review.
% The framework supports validation of kernel optimizations, verified translation
% across ISAs, and cross-checking of existing implementations. The planned main
% study uses SALTyRN's 35 published final optimized RVV kernels. Additional studies
% target AVX-to-RVV/HVX ports, XNNPACK implementation comparisons, and a composed
% SIMD program. All proof counts, repairs, cost measurements, and speedups remain
% \todo{pending evaluation}; none is asserted as a completed contribution.
